\chapter{Height barriers and entire graphs of arbitrary growth}\label{sec:graphical-height-existence}
Shrinking spheres give height control without differentiating the initial graph. We use the avoidance argument of \cite{EH91}, Section 5, pp.~567--568, and keep the subsequent parabolic existence theorem explicitly external. The assumptions for bounded-slope comparison are not silently transferred to arbitrary growth.

\begin{lemma}[Avoidance from a locally continuous initial graph]\label{lem:entire-graph-initial-trace-avoidance}
Let an entire graphical mean curvature flow be smooth for $0<t\leq T$ and converge locally uniformly as $t\downarrow0$ to a continuous graph. Let $N_t$ be a compact smooth hypersurface flow on $[0,T]$, whose initial image is disjoint from that graph. Then the two positive-time images remain disjoint, and their distance is at least their strictly positive initial distance.
\end{lemma}
\begin{proof}
The initial distance is attained and positive: a compact hypersurface and a closed entire graph are disjoint, and graph points with large base coordinate cannot minimize their distance. Over the whole time interval the compact flow has bounded image. Testing distances against a fixed base point bounds the minimum distance above uniformly, since the graph heights there extend continuously to time zero. Thus all minimizing graph base points lie in one compact ball. Locally uniform initial convergence makes the minimum distances converge to the initial positive distance as $t\downarrow0$. On $[\epsilon,T]$, Lemma \ref{lem:compact-flow-entire-graph-avoidance} makes the distance nondecreasing. Let $\epsilon\downarrow0$ to obtain the asserted lower bound. No tangent plane or curvature for the initial graph was used.
\end{proof}

\begin{proposition}[Local heights are controlled by a larger initial ball]\label{prop:entire-graph-local-height-barrier}
In the preceding smooth positive-time, locally continuous initial-trace class, fix a base point $y$, $r\geq0$, $T>0$, and $R$ with $R^2>r^2+2nT$. If $|u(x,0)|\leq M$ for $|x-y|\leq R$, then for $|x-y|\leq r$ and $0\leq t\leq T$,
\[
 |u(x,t)|\leq M+R-\sqrt{R^2-2nt-|x-y|^2}.
\]
In particular this is a finite bound determined by the initial graph on a compact base ball. This is a quantitative form of the sphere barriers in \cite{EH91}, proof of Theorem 5.1.
\end{proposition}
\begin{proof}
Place a sphere of radius $R$ with center $(y,M+R+\epsilon)$, for $\epsilon>0$. Its entire ball lies strictly above the initial graph over its base projection, because its lowest height is $M+\epsilon$. Outside that projection the graph cannot meet the ball. The sphere shrinks to radius $\sqrt{R^2-2nt}$ and stays smooth through $T$. Lemma \ref{lem:entire-graph-initial-trace-avoidance} prevents contact. At each fixed base point in $B_r(y)$ its lower sheet exists throughout the interval, and the graph begins below that sheet. Continuity prevents crossing. The resulting upper bound is the displayed expression with an added $\epsilon$; let $\epsilon\downarrow0$. Reflect the construction for the lower bound. A bound on the initial height outside $B_R(y)$ is unnecessary because the initial comparison sphere has that bounded base projection.
\end{proof}

\begin{theorem}[A quantitative half-Hölder bound from Lipschitz initial data]\label{thm:graph-lipschitz-half-holder-time}
Let a smooth positive-time entire graphical flow converge locally uniformly to a globally $L$-Lipschitz initial function $u_0$, with $L\geq0$. Then
\[
 |u(x,t)-u_0(x)|\leq L\sqrt{2nt}
 \qquad(x\in\mathbb R^n,\ t>0)
\]
at every time of smooth existence. A slope or curvature bound for the later flow is not an additional hypothesis. The sphere-barrier method is \cite{EH91}, Section 5; the optimized constant is derived here.
\end{theorem}
\begin{proof}
Fix $x=y$ and $t>0$. For $R>\sqrt{2nt}$, place a sphere of radius $R$ with center at height $u_0(y)+R\sqrt{1+L^2}+\epsilon$ over $y$. Its lower sheet lies above the whole initial graph in its projected ball, since for $0\leq r\leq R$,
\[
 Lr+\sqrt{R^2-r^2}\leq R\sqrt{1+L^2}.
\]
This is Cauchy--Schwarz applied to $(L,1)$ and $(r,\sqrt{R^2-r^2})$. The initial Lipschitz bound gives $u_0(x)\leq u_0(y)+L|x-y|$, so the vertical gap is at least $\epsilon$. Avoidance and continuity below the lower sheet, as in Proposition \ref{prop:entire-graph-local-height-barrier}, imply
\[
 u(y,t)-u_0(y)\leq R\sqrt{1+L^2}-\sqrt{R^2-2nt}
\]
after letting $\epsilon\downarrow0$. For $L>0$, choose $R=\sqrt{2nt}\sqrt{1+L^2}/L$; substitution makes the right side $L\sqrt{2nt}$. For $L=0$, let $R\to\infty$ to obtain zero. A reflected sphere gives the lower estimate. Since $y$ was arbitrary, the conclusion is uniform in space although initial heights themselves may be unbounded.
\end{proof}

\begin{proposition}[Uniform contraction and a time modulus for bounded-slope solutions]\label{prop:graphical-uniform-contraction}
For the globally Lipschitz-data solutions of Corollary \ref{cor:globally-lipschitz-graph-wellposedness},
\[
 \sup_x|u(x,t)-\widetilde u(x,t)|
 \leq\sup_x|u(x,0)-\widetilde u(x,0)|
\]
whenever the right side is finite. If the initial Lipschitz constant is at most $L$, then
\[
 \sup_x|u(x,t)-u(x,s)|\leq L\sqrt{2n(t-s)}\qquad(0\leq s<t).
\]
\end{proposition}
\begin{proof}
If the initial difference is bounded by $\delta$, the functions $\widetilde u\pm\delta$ are also graphical solutions. Theorem \ref{thm:entire-bounded-slope-comparison} traps $u$ between them, proving contraction. The existence theorem gives $|Du|\leq L$ at all positive times. Start Theorem \ref{thm:graph-lipschitz-half-holder-time} at time $s$ with the function $u(\cdot,s)$; the same Lipschitz constant applies. This proves the second estimate, with $s=0$ directly covered by that theorem. The height modulus uses compact sphere barriers; it is not obtained by integrating a nonexistent uniform initial curvature bound.
\end{proof}

\begin{theorem}[Ecker--Huisken existence for entire locally Lipschitz graphs]\label{thm:ecker-huisken-arbitrary-growth-existence}
Every locally Lipschitz function $u_0:\mathbb R^n\to\mathbb R$, $n\geq1$, is the initial value of an entire graphical mean curvature flow that is smooth for all $t>0$ and converges locally uniformly to $u_0$ at time zero. No global slope, height-growth or curvature bound is required. This is \cite{EH91}, Theorem 5.1 and Remark 5.2(i), printed p.~567, with its exhaustion and approximation argument on p.~568.

The existence proof remains external: it uses graphical boundary-value existence, height-dependent local gradient estimates, higher derivative estimates and compact exhaustion. The local height and curvature estimates here and the height-dependent gradient proof in Theorem \ref{thm:graphical-height-gradient-estimate} explain parts of that argument but do not replace its remaining steps. Uniqueness in the unrestricted unbounded-slope class is not asserted by this statement.
\end{theorem}

\begin{example}[A rapidly growing graph and the divergent heat-kernel integral]\label{ex:graph-flow-superexponential-data}
The smooth initial graph $u_0(x)=e^{|x|^4}$ admits an all-time positive-time smooth entire graphical mean curvature flow by Theorem \ref{thm:ecker-huisken-arbitrary-growth-existence}. In contrast, for every $t>0$ and fixed $x$, the Euclidean heat-kernel integral
\[
 (4\pi t)^{-n/2}\int_{\mathbb R^n}e^{-|x-y|^2/(4t)}e^{|y|^4}\,dy
\]
is infinite. This comparison concerns the heat-kernel construction, not a classification of all possible sign-changing heat-equation solutions without growth restrictions.
\end{example}
\begin{proof}
The initial function is smooth and hence locally Lipschitz, so it meets the cited graphical existence theorem despite its unbounded slope and rapid growth. For large $|y|$ with $x,t$ fixed, the exponent $|y|^4-|x-y|^2/(4t)$ is at least $|y|^4/2$. The integral over the exterior of that ball therefore diverges. The graph theorem and the heat-kernel computation make different assertions, and neither is deduced from a formal linearization at a graph with unbounded gradient. Local finite-height bounds for the nonlinear flow are consistent with Proposition \ref{prop:entire-graph-local-height-barrier}, which only asks for initial heights on a finite larger ball.
\end{proof}
